\noindent\hspace*{1.2em}\kbvlabel{ಉದಾ}\hspace{0.6em}ಸೂರ್ಯನು ಚೆನ್ನಾಗಿ ಪ್ರಕಾಶಿಸುತ್ತಿದ್ದಾನೆ. (ಪುಲ್ಲಿಂಗ)\par

\indent ಬೆಳಿಗ್ಗೆ ಆರು ಗಂಟೆಗೆ ಸೂರ್ಯೋದಯವಾಯಿತು. (ನಪುಂಸಕ ಲಿಂಗ)\par

\begin{kbvrule}
\noindent\kbvbullet\hspace{0.5em}ಕೆಲವು ಶಬ್ದಗಳನ್ನು ಸ್ತ್ರೀಲಿಂಗದಲ್ಲಿ ಹಾಗೂ ನಪುಂಸಕ ಲಿಂಗದಲ್ಲಿ ಬಳಸುವ ವಾಡಿಕೆಯುಂಟು. ಅಂತಹ ಶಬ್ದಗಳು ಸ್ತ್ರೀನಪುಂಸಕ ಲಿಂಗಕ್ಕೆ ಸೇರಿದವುಗಳೆಂದು ಹೇಳುತ್ತೇವೆ.
\end{kbvrule}

\noindent\hspace*{1.2em}\kbvlabel{ಉದಾ}\hspace{0.6em}ಅಣ್ಣಮ್ಮದೇವಿ ವರ ಕೊಟ್ಟಳು. (ಸ್ತ್ರೀಲಿಂಗ)\par

\begin{kbvrule}
\noindent\kbvbullet\hspace{0.5em}ಅಣ್ಣಮ್ಮದೇವಿ ವರ ಕೊಟ್ಟಿತು. (ನಪುಂಸಕ ಲಿಂಗ)\par
\noindent\kbvbullet\hspace{0.5em}ಕೆಲವು ಶಬ್ದಗಳನ್ನು ಮೂರು ಲಿಂಗಗಳಲ್ಲೂ ಬಳಸುವ ಸಾಧ್ಯತೆ ಯುಂಟು. ಅಂತಹ ಶಬ್ದಗಳು ವಿಶೇಷ್ಯಾಧೀನ ಲಿಂಗಕ್ಕೆ ಅಥವಾ ವಾಚ್ಯ ಲಿಂಗಕ್ಕೆ ಸೇರಿದವುಗಳೆಂದು ಹೇಳುತ್ತೇವೆ.
\end{kbvrule}

\noindent\hspace*{1.2em}\kbvlabel{ಉದಾ}\hspace{0.6em}. ಅವನು ಜೀತದ ಆಳು. (ಪು.ಲಿಂ.) ಆ\par

\indent ಮನೆಗೆಲಸದ ಆಳು ಈ ದಿನ ಬಂದಿದ್ದಳು. (ಸ್ತ್ರೀ ಲಿಂ.) ಊರಿನಿಂದ ಆಳು ಬಂದಿತ್ತು. (ನ.ಲಿಂ.) ನಾನು ನಿನ್ನ ತಂದೆ. (ಪು.ಲಿಂ.) ನಾನು ನಿನ್ನ ತಾಯಿ. (ಸ್ತ್ರೀ.ಲಿಂ.) ನಾನು ನಿಮ್ಮ ಮಗು. (ನ.ಲಿಂ.)\par

\indent ವಿ. ಸೂ. : ಪ್ರಧಾನಿ, ಮಂತ್ರಿ, ವಿಜ್ಞಾನಿ, ಪ್ರಜೆ, ಕಾರ್ಯದರ್ಶಿ ಈ ಮೊದಲಾದ ಶಬ್ದಗಳು ವಾಚ್ಯಲಿಂಗ ಶಬ್ದಗಳಾಗಿರುವುವು.\par

\begin{kbvrule}
\noindent\kbvbullet\hspace{0.5em}ಕೆಲವು ಶಬ್ದಗಳನ್ನು ನಪುಂಸಕ ಲಿಂಗದಲ್ಲಿ ಮಾತ್ರ ಬಳಸುವುದು - ವಾಡಿಕೆ :: ಅವು ಪುಲ್ಲಿಂಗವಾದರೂ ಆಗಿರಬಹುದು. ಸ್ತ್ರೀಲಿಂಗವಾದರೂ ಆಗಿರಬಹುದು. ಅಂತಹ ಶಬ್ದಗಳು ನಿತ್ಯನಪುಂಸಕ ಲಿಂಗಕ್ಕೆ ಸೇರಿದವು - ಗಳೆಂದು ಹೇಳುತ್ತೇವೆ.
\end{kbvrule}

\noindent\hspace*{1.2em}\kbvlabel{ಉದಾ}\hspace{0.6em}ಮಗು, ಕೂಸು, ಶಿಶು, ಕಂದ, ಹಸುಳೆ, ದಂಡು, ಸೇನೆ, ದಳ, ಜನ ಇತ್ಯಾದಿ. \par
